\documentclass[12pt]{article}
\usepackage[margin=1in]{geometry}
\usepackage{parskip}
\usepackage{hyperref}

\hypersetup{
    colorlinks=true,
    linkcolor=blue,
    urlcolor=blue
}

\title{\textbf{Snapshot Objectives} \\ EagleNav -- Campus Navigation Using Augmented Reality}
\author{Team 2 \\
Anna Guerrero-Leon, Jose Mateo Ayala, Justin Ho, Min Park \\
California State University, Los Angeles}
\date{May 15, 2026}

\begin{document}

\maketitle
\newpage

\tableofcontents
\newpage

\section{Start Objective}

The goal of Snapshot 1 is to establish the foundation of the EagleNav project. The team will set up the GitHub repository, configure the Jira board with the first sprint, and confirm the core tech stack: Flutter for the mobile app, Valhalla for routing, and OpenStreetMap for campus map data. Initial drafts of the SDD and SRS will be started, and a high-level Workflow Diagram will be created to outline the overall system architecture.

\section{1st Checkpoint Objective}

In Snapshot 1, the team established the project foundation by setting up the repository, Jira board, and beginning the SDD and SRS. For Snapshot 2, the focus shifts to documenting the core navigation system. This includes the layered system architecture, the route generation flow, and the GPS-based guidance loop that drives turn-by-turn navigation. The SRS, Workflow Diagram, and Design Spec will all be updated to reflect these additions.

\section{2nd Checkpoint Objective}

In Snapshot 2, the team documented the routing and guidance systems. For Snapshot 3, the focus is on the voice guidance and accessibility layer, which is the feature that sets EagleNav apart from standard navigation apps. This includes the Orientation Coach, Instruction Enhancer, and Landmark Service, all of which provide body-relative spoken directions to support visually impaired users. The SDD, SRS, and all supporting documents will be updated to reflect these new features.

\section{Due Date Checkpoint}

In Snapshot 3, the team documented the voice guidance and accessibility layer including orientation coaching, instruction enhancement, and the landmark service. For Snapshot 4, the focus is on finalizing the object detection and perception layer, which allows the app to detect real-time obstacles using the device camera. All documents will receive their final updates, and future work will be noted, including indoor navigation, sensor fusion, and stride length calibration.

\end{document}